%=====================================================================
\chapter{Presenting and Defending Your Research}\label{ch:defending}
%=====================================================================
\locator{6}{Part~6 --- Producing and Defending Research}

\begin{outcomes}
By the end of this chapter you will be able to:
\begin{tight}
  \item \textbf{Describe} each examination you will face and what it assesses.
  \item \textbf{Prepare} for the comprehensive examination against its actual
    paper pattern.
  \item \textbf{Build} a defence presentation that survives interruption.
  \item \textbf{Answer} the six question types examiners use.
  \item \textbf{Interpret} the four verdicts and what follows from each.
\end{tight}
\end{outcomes}

\begin{prereq}
The whole course. This chapter is where you are asked to account for every
decision Parts~I--V asked you to make deliberately.
\end{prereq}

\begin{keyterms}
comprehensive examination $\bullet$ synopsis defence $\bullet$ viva voce
$\bullet$ open defence $\bullet$ internal and external examiner $\bullet$
minor amendments $\bullet$ deferred for resubmission $\bullet$ mock defence
\end{keyterms}

\section{Three Examinations, Different Purposes}

\begin{center}
\small
\begin{longtable}{@{}L{3.0cm}L{5.2cm}L{6.0cm}@{}}
\toprule
\textbf{Examination} & \textbf{Assesses} & \textbf{Rules} \\
\midrule
\endfirsthead
\toprule
\textbf{Examination} & \textbf{Assesses} & \textbf{Rules} \\
\midrule
\endhead
\bottomrule
\endfoot
Comprehensive (PhD only)
  & Doctoral-level mastery of your specialisation --- not your own research
  & \S11. Within 60 days of coursework at CGPA 3.0; written and oral; 65\% to
    pass each part; internal examiners are those who taught you, plus an
    external appointed by the Dean; one retake within 30 days, then dropped \\
Synopsis defence
  & Whether the proposed research should proceed
  & PhD \S13 before the DRC, plus external referee. MS \S8 before the DRC
    (\chref{synopsis}) \\
Viva voce / open defence
  & Whether the completed research merits the degree
  & PhD \S33: open defence, publicly announced, one external national examiner
    present, after positive foreign reports. MS \S48: external examiner, open
    defence preferred \\
\end{longtable}
\end{center}

\begin{regbox}[The comprehensive examination paper pattern --- MPhil/MS \S12(v)]
Where a comprehensive examination is held, the pattern is specified: True/False
20 marks, Multiple Choice 20 marks, Short Questions 20 marks, and an Oral
Examination separately worth 40 marks. The passing score for both the oral and
the written parts is 60\% at MS level, and \reg{PhD Regs 2024}{\S11(iii)}
sets it at 65\% for PhD.

\medskip
Note what this implies: the oral is worth as much as the entire written paper.
Prepare for it as seriously.
\end{regbox}

\begin{conceptbox}[What the comprehensive examination is really testing]
\S11(ii) says it ``should be based on recapitulation of the conceptual
knowledge of the discipline to which the student is admitted. Questions may be
included therein to test the overall grasp of the relevant field achieved by
the student.''

\medskip
It is breadth, not depth in your own topic. The papers you took are the
syllabus, and methodology questions come from this course. Revise from the
course portfolio (Appendix~\ref{app:schedule}) rather than from your thesis
reading, which is far too narrow for this purpose.
\end{conceptbox}

\section{The Defence Presentation}

\begin{center}
\small
\begin{tabular}{@{}L{1.6cm}L{3.6cm}L{8.7cm}@{}}
\toprule
\textbf{Slides} & \textbf{Content} & \textbf{Note} \\
\midrule
1     & Title, you, supervisor, date & \\
2     & The problem                  & In the world, not in the literature \\
3     & The gap                      & With the evidence for it \\
4     & Research questions           & Numbered; they will be referred to \\
5--6  & Method                       & Design and why, not procedure detail \\
7--10 & Results                      & One finding per slide, with uncertainty \\
11    & Threats                      & \textbf{Include it.} Naming your own
                                       limits first is the strongest move
                                       available (\chref{validity}) \\
12    & Contribution                 & As a list (\chref{writing}) \\
13    & Future work                  & Two or three specific studies \\
\midrule
Backup & Twenty or more             & Every table, every alternative analysis,
                                      every ablation. This is where a defence
                                      is won \\
\bottomrule
\end{tabular}
\end{center}

\begin{alertbox}[Backup slides are the preparation that shows]
When an examiner asks ``did you check whether the effect holds for the
subgroup?'' there are two possible answers. ``No, but that would be interesting
future work'' is survivable. Pulling up slide B14, which contains exactly that
analysis, ends the line of questioning and changes the room.

\medskip
Build backup slides for every question you can imagine, and especially for the
ones you hope nobody asks. Preparing them also forces you to actually run the
analyses you have been avoiding --- which is the real benefit.
\end{alertbox}

\begin{center}
\small
\begin{tabular}{@{}L{3.4cm}L{10.9cm}@{}}
\toprule
Design for interruption & You may not reach slide 13. Front-load: the
                          contribution should be stated by slide 4 as well as
                          slide 12 \\
One idea per slide      & If it needs two minutes of explanation, it is two
                          slides \\
Figures, not tables     & A table of twenty numbers is unreadable projected.
                          Plot it; keep the table in backup \\
No slide read aloud     & If you read it, they read ahead and stop listening \\
Rehearse aloud, timed   & Silently in your head is not rehearsal, and the
                          difference is roughly forty per cent on timing \\
\bottomrule
\end{tabular}
\end{center}

\section{Questions}

\begin{center}
\small
\begin{longtable}{@{}L{3.0cm}L{4.8cm}L{6.4cm}@{}}
\toprule
\textbf{Type} & \textbf{Example} & \textbf{How to answer} \\
\midrule
\endfirsthead
\toprule
\textbf{Type} & \textbf{Example} & \textbf{How to answer} \\
\midrule
\endhead
\bottomrule
\endfoot
Clarification
  & ``What exactly is the unit of analysis?''
  & Answer directly, in one sentence. Do not elaborate \\
Justification
  & ``Why a case study rather than a survey?''
  & Name the alternative, the criterion, and the constraint that decided it.
    This is portfolio piece P3, which is why you wrote it \\
Challenge
  & ``Your sample cannot support that conclusion.''
  & Restate the objection accurately first. Then either concede and scope the
    claim, or give the evidence. Never argue past it \\
Extension
  & ``What would happen with distributed teams?''
  & Answer as a researcher: what you would predict, and what study would test
    it \\
Knowledge
  & ``What is the difference between reliability and validity?''
  & They are checking foundations. Answer plainly and completely; this is
    \chref{measurement} \\
Hostile or confused
  & A question resting on a misreading
  & Clarify gently --- ``I may not have made this clear in the thesis'' --- then
    answer the question they meant to ask \\
\end{longtable}
\end{center}

\begin{conceptbox}[Four habits worth rehearsing]
\begin{tightnum}
  \item \textbf{Pause before answering.} Two seconds of thought reads as
    consideration, not hesitation. Answering instantly reads as rehearsed.
  \item \textbf{Answer the question asked.} Not the adjacent one you prepared
    for. Examiners notice the substitution immediately.
  \item \textbf{Say ``I don't know'' when you do not.} Then add what you would
    do to find out. This is a good answer; bluffing is the only bad one, and
    it is always detected.
  \item \textbf{Concede what is true.} An examiner who lands a real hit and is
    met with agreement --- ``Yes, that is a genuine limitation, and here is
    what it does and does not affect'' --- generally moves on. One met with
    defensiveness will press.
\end{tightnum}
\end{conceptbox}

\begin{alertbox}[The two questions that come up almost always]
\textbf{``What is your contribution?''} Have one sentence ready
(\chref{contribution}, \chref{writing}). Fumbling this is the single worst
moment available in a defence, and it is entirely preventable.

\medskip
\textbf{``What would you do differently?''} This is not a trap. It tests
whether you have judgement about your own work. Have a real answer --- a design
choice you would change and why --- and avoid both the false modesty of
``nothing'' and a list so long it suggests the work was poorly planned.
\end{alertbox}

\section{The Verdicts}

\begin{center}
\small
\begin{tabular}{@{}L{4.0cm}L{5.0cm}L{5.3cm}@{}}
\toprule
\textbf{Verdict} & \textbf{PhD \S33(vii)} & \textbf{MS \S48(xi)} \\
\midrule
Pass
  & Unconditional
  & Unconditional \\
Pass with minor amendments
  & Corrections, then done
  & Corrections, then done \\
Deferred for resubmission and re-defence
  & Re-defence rearranged within \textbf{180 days} after corrections,
    acknowledged by the supervisor --- otherwise ``fail''
  & Resubmission in revised form within \textbf{six months} of intimation from
    DASR; prescribed fee payable \\
Fail
  & ``In exceptional circumstances and for the reasons to be recorded by the
    examiners''
  & Recorded by the examiner \\
\bottomrule
\end{tabular}
\end{center}

\begin{conceptbox}[Deferred is not failed, and the clock is real]
Most deferrals are recoverable. What ends them badly is the deadline: the PhD
regulation is explicit that if the corrected version is not acknowledged within
180 days, the candidate ``will be considered fail''.

\medskip
So on the day you receive a deferral: get the required changes in writing, agree
a schedule with your supervisor immediately, and start. The temptation to put
the thesis aside for a month after a difficult defence is understandable and is
the single most expensive month available to you.
\end{conceptbox}

\section{Preparing}

\begin{center}
\small
\begin{tabular}{@{}L{3.4cm}L{10.9cm}@{}}
\toprule
Re-read the thesis  & Completely, a week before. You wrote Chapter 3 eighteen
                      months ago and will be asked about it \\
List every decision & Every design and analysis choice, with the alternative
                      and the reason. This list \emph{is} the defence, and P3
                      and P5 already contain most of it \\
Prepare the weak points & Write out the three hardest questions and your
                      answers. Examiners find the same weak points you know
                      about \\
Mock defence        & With colleagues instructed to be difficult. Non-negotiable \\
Know the examiners  & Read a recent paper by each. It tells you what they care
                      about, and it is a courtesy \\
Logistics           & Room, projector, adapter, backup of the slides on a
                      second device and in email \\
Dues and clearance  & \S33(i): the viva is conducted ``only if he/she has
                      cleared all dues and has paid full prescribed examination
                      fees''. Check weeks in advance \\
\bottomrule
\end{tabular}
\end{center}

\begin{traditions}{the defence}
\tradEMP{Expect challenges to power, measurement and causal warrant. Have the
interval, the reliability coefficient and the identification argument ready.}
\tradDSR{Expect ``what is the contribution beyond the artefact?''. Have the
design knowledge sentence (\chref{designscience}).}
\tradQUAL{Expect questions about sample size and generalisability framed in
quantitative terms. Answer with saturation evidence and analytic
generalisation, and do so without conceding the frame (\chref{qda}).}
\tradFORM{Expect the assumptions to be attacked rather than the proof. Know
which assumptions are essential and which are convenience.}
\end{traditions}

\begin{lab}[Run a real mock defence]
\begin{tightnum}
  \item Build the 13-slide deck plus at least fifteen backup slides.
  \item Present to three colleagues who have read the work, with instructions
    to interrupt and to press when an answer is thin.
  \item Have someone log every question and how well you answered it, on a
    three-point scale.
  \item Build a backup slide for every question you scored below the top mark.
  \item Repeat with a different audience.
  \item Write your one-sentence contribution on a card and carry it.
\end{tightnum}
\end{lab}

\begin{checkpoint}
\begin{tightnum}
  \item What does the comprehensive examination assess, and why is revising
    from your thesis reading the wrong preparation?
  \item Why should the contribution appear by slide 4 rather than only at the
    end?
  \item An examiner makes a valid criticism you cannot answer. What do you say?
  \item How long do you have after a deferral, and what happens if you miss it?
\end{tightnum}
\end{checkpoint}

\begin{chaptersummary}
\begin{tight}
  \item Three examinations: comprehensive (PhD, breadth, 65\% written and oral,
    one retake), synopsis defence (should this proceed), viva (does this merit
    the degree).
  \item The comprehensive oral can carry as much weight as the whole written
    paper. Prepare it accordingly.
  \item Thirteen slides, contribution by slide 4, threats included, and a large
    set of backup slides. Backup slides are where defences are won.
  \item Six question types. Pause, answer what was asked, say ``I don't know''
    when true, and concede what is true.
  \item Have the contribution sentence and the ``what would you do
    differently'' answer ready.
  \item Deferred is recoverable but the clock is real --- 180 days for PhD, six
    months for MS. Start immediately.
  \item Mock defence with hostile colleagues, twice.
\end{tight}
\end{chaptersummary}

\begin{reviewq}
  \item The PhD open defence is publicly announced and anyone may attend and
    question (\S33(v)--(vi)). What does public defence add over a closed
    examination, and what does it cost the candidate?
  \item Qualitative theses are routinely challenged on sample size by examiners
    applying quantitative criteria. How should a candidate respond in the
    moment --- and whose responsibility is it that the question arises?
  \item \S11 allows one retake of the comprehensive examination, after which
    the scholar is dropped. Is a single retake the right policy? Consider both
    standards and the sunk investment of scholar and department.
  \item Design a mock defence programme for your department. Who participates,
    at what stage, and how would you know whether it improved outcomes?
\end{reviewq}
